%% Hand-authored circuitikz figure -- NOT generated by maestro2tex, placed by it via
%% the `order` list in configs/anatop_tb_tia_amp_ab_corners.json. Topology transcribed
%% from the class-AB TIA amp bench netlists (Interactive.14, archived under
%% simarchive/tb_anatop_tia_amp_ab/tia_ab_2026-08-16/). Supplies, the bias rails and
%% En are omitted -- identical in every bench and carrying no measurement.
\begin{figure}[htbp]
  \centering
  \resizebox{0.56\linewidth}{!}{%
\ctikzset{bipoles/length=1.0cm, resistors/scale=0.85, capacitors/scale=0.85, sources/scale=0.85}
\begin{circuitikz}[
    every node/.append style={font=\footnotesize},
    nlab/.style={font=\scriptsize, inner sep=1.5pt},
    opt/.style={font=\scriptsize, text=black!55, align=center, inner sep=2pt},
    railnode/.style={circle, fill=black, inner sep=1.05pt},
    prb/.style={draw, circle, minimum size=9pt, inner sep=0pt},
  ]
  \node[op amp] (B) at (0,0) {};
  \node at (B.center) {\footnotesize A2};
  %% Input pins leave horizontally; every branch that meets them is placed at the
  %% pin height with (x,0 |- pin) so no run is ever diagonal.
  \coordinate (Bp) at (B.+);
  \coordinate (Bm) at (B.-);
  \coordinate (VCM) at (-1.6,0 |- Bp);
  \coordinate (WE)  at (-3.0,0 |- Bm);
  \coordinate (ZC)  at (-4.6,0 |- Bm);
  \draw (B.+) -- (VCM);
  \draw (VCM) to[V, l_=$V_{\mathrm{CM}}$, invert] (-1.6,-2.8) node[ground]{};
  \draw (B.-) -- (WE);
  \node[railnode] at (WE){};
  \node[nlab, anchor=south east, xshift=-1pt, yshift=2pt] at (WE){\texttt{WE}};
  \draw (WE) to[I, l_=$i_{\mathrm{we}}$] (-3.0,-2.8) node[ground]{};
  \draw (WE) -- (ZC);
  \draw (ZC) to[generic, l_=$Z_{\mathrm{cell}}$] (-4.6,-2.4);
  \node[nlab, below] at (-4.6,-2.45) {$v_{\mathrm{cm}}$};
  \draw (WE) -- (-3.0,2.2);
  \node[prb] (PB) at (-2.0,2.2) {};
  \draw (-3.0,2.2) -- (PB.west);
  \node[nlab, above=2pt] at (PB.north) {probe};
  \draw (PB.east) -- (-1.2,2.2);
  \draw (-1.2,2.2) to[R, l=$R_f$] (1.6,2.2);
  \draw (-1.2,2.2) -- (-1.2,3.5);
  \draw (-1.2,3.5) to[C, l=$C_f$] (1.6,3.5);
  \draw (1.6,3.5) -- (1.6,0);
  \draw (B.out) -- (1.6,0);
  \node[railnode] at (1.6,0){};
  \node[nlab, below right] at (1.7,-0.05){$v_{\mathrm{out}}$};
  \draw (1.6,0) -- (3.0,0);
  \draw (3.0,0) to[C, l=$C_L$] (3.0,-2.8) node[ground]{};
\end{circuitikz}
  }
  \caption[{The closed transimpedance loop bench.}]{The closed transimpedance loop bench, the only one in this section that closes feedback around the electrode. $R_f \parallel C_f$ returns the output to the working electrode through a current probe, $i_{\mathrm{we}}$ is the test source, and $Z_{\mathrm{cell}}$ is the Randles cell of Figure~\ref{fig:randles-cell}, returned to $v_{\mathrm{cm}}$. The loop gain, phase margin and crossover are measured through the probe; the DC accuracy, closed-loop transfer and noise figures come from the same wiring.}
  \label{fig:tb-tiafb}
\end{figure}
